\documentclass[10pt,aspectratio=1610]{beamer}

% ================================================================================
% Load SMU Theme
% ================================================================================
\usepackage{smu}

% ================================================================================
% Bibliography Configuration (User-defined)
% ================================================================================
\usepackage[
    style=alphabetic,
    block=space,
    maxbibnames=99,
    maxalphanames=3,
    minalphanames=3
]{biblatex}
\addbibresource{ref.bib}

% Custom citation label template
\DeclareLabelalphaTemplate{
  \labelelement{
    \field[final]{shorthand}
    \field{label}
    \field[strwidth=3,strside=left,ifnames=1]{labelname}
    \field[strwidth=1,strside=left,ifnames=2]{labelname}
    \field[strwidth=1,strside=left,ifnames=3]{labelname}
    \field[strwidth=1,strside=left]{labelname}
  }
  \labelelement{
    \literal{+}
  }
  \labelelement{
    \field[strwidth=2,strside=right]{year}
  }
}

\renewcommand*{\labelalphaothers}{\ifnumgreater{\value{listtotal}}{3}{+}{}}
\renewcommand*{\newunitpunct}{\addcomma\space}

% Field formatting
\DeclareFieldFormat*{title}{#1\addperiod}
\DeclareFieldFormat*{citetitle}{#1}
\DeclareFieldFormat{booktitle}{\textit{#1}}
\DeclareFieldFormat{journaltitle}{\textit{#1}}

% Bibliography spacing for beamer
\setlength{\bibitemsep}{0.3ex}
\setlength{\bibnamesep}{0.3ex}
\setlength{\bibinitsep}{0ex}
\renewcommand*{\bibfont}{\tiny}

% ================================================================================
% Additional Packages
% ================================================================================
\usepackage{amsmath,amssymb}
\usepackage{booktabs}
\usepackage{lipsum}

% ================================================================================
% Presentation Metadata
% ================================================================================
\title{SMU Beamer Theme}
\subtitle{Complete Feature Demonstration}
\author{Your Name}
\institute[SMU]{School of Computing and Information Systems\\Singapore Management University}
\date{\today}

% ================================================================================
% Document Body
% ================================================================================
\begin{document}

% ================================================================================
% Title Page
% ================================================================================
\maketitle

% ================================================================================
% Table of Contents
% ================================================================================
\begin{frame}[allowframebreaks]{Overview}
    \tableofcontents
\end{frame}

% ================================================================================
\section{Basic Settings}
% ================================================================================

\subsection{Introduction}

\begin{frame}{About This Theme}
    The SMU Beamer Theme provides a LaTeX beamer template based on SMU's official PowerPoint template, with enhanced features for academic presentations:
    \begin{itemize}
        \item Official SMU branding and colors
        \item Code highlighting with minted
        \item Custom block components
        \item Flexible layout options
    \end{itemize}
\end{frame}

\begin{frame}{Key Features}
    \begin{twocolumns}
        \leftcol{
            \textbf{Design \& Layout:}
            \begin{itemize}
                \item Professional SMU branding
                \item Progress bar footer
                \item Section transition pages
                \item Multi-column layouts
            \end{itemize}
            
            \vspace{0.5em}
            
            \textbf{Text \& Colors:}
            \begin{itemize}
                \item Custom color commands
                \item Text highlighting
                \item Bold color variants
            \end{itemize}
        }
        \rightcol{
            \textbf{Code Highlighting:}
            \begin{itemize}
                \item Minted integration
                \item 10+ programming languages
                \item 49 built-in themes
                \item Inline code support
            \end{itemize}
            
            \vspace{0.5em}
            
            \textbf{Custom Blocks:}
            \begin{itemize}
                \item Standard blocks
                \item Alert blocks
                \item Mathematical boxes (Theorem, etc.)
                \item Tagged blocks
            \end{itemize}
        }
    \end{twocolumns}
\end{frame}

% ================================================================================
\subsection{Text Formatting}
% ================================================================================

\begin{frame}{Basic Text Formatting}
    \textbf{Font Styles:}
    \begin{itemize}
        \item Regular text in Calibri font
        \item \textbf{Bold text appears in SMU blue}
        \item \textit{Italic text for emphasis}
        \item \redbf{Important warning} using red bold
        \item \bluebf{Key concept} using blue bold
    \end{itemize}
    
    \vspace{1em}
    
    \textbf{Inline Code Style:}
    \begin{itemize}
        \item Use \smucode{inline code} for technical terms
        \item Function names like \smucode{printf()} or \smucode{main()}
        \item Styles can be customized in \smucode{smucore/smucode.def}
    \end{itemize}
\end{frame}

\begin{frame}{Custom Colors}
    \textbf{SMU Brand Colors:}
    \begin{itemize}
        \item \smucolor{smublue}{SMU Blue - Primary color}
        \item \smucolor{smulightblue}{SMU Light Blue - Secondary color}
        \item \smucolor{smugold}{SMU Gold - Accent color}
        \item \smucolor{smugreen}{SMU Green - Additional accent}
    \end{itemize}
    
    \vspace{1em}
    
    \textbf{Additional Accent Colors:}
    \begin{itemize}
        \item \smucolor{scisyellow}{SCIS Yellow - School of Computing}
        \item \smucolor{alertred}{Alert Red - For warnings}
        \item \smucolor{mutedblue}{Muted Blue - Alternative blue}
        \item \smucolor{mutedgreen}{Muted Green - Alternative green}
    \end{itemize}
\end{frame}

\begin{frame}{Mathematical Formulas}
    \textbf{Inline Mathematics:}
    
    Einstein's famous equation $E = mc^2$, Euler's number $e \approx 2.71828$, and pi $\pi \approx 3.14159$.
    
    \vspace{1em}
    
    \textbf{Display Equations:}
    \[ \int_0^\infty e^{-x^2} dx = \frac{\sqrt{\pi}}{2} \]
    
    \[ \nabla \times \mathbf{E} = -\frac{\partial \mathbf{B}}{\partial t} \]
    
    \vspace{0.5em}
    
    \textbf{Aligned Equations:}
    \begin{align*}
        f(x) &= x^2 + 2x + 1 \\
             &= (x + 1)^2 \\
             &= (x + 1)(x + 1)
    \end{align*}
\end{frame}

\begin{frame}{Text \& Mathematical Highlighting}
    \textbf{Basic Highlighting:}
    
    Normal text with \shadedtext{default highlighted content} using SMU blue.
    
    \vspace{0.5em}
    
    \textbf{Custom Colors:}
    
    Different colors: \shadedtext[alertred]{red highlight}, 
    \shadedtext[smugreen]{green highlight}, 
    \shadedtext[scisyellow]{yellow highlight}, and 
    \shadedtext[mutedblue]{blue highlight}.
    
    \vspace{1em}
    
    \textbf{Math Formula Highlighting:}
    
    Inline math: \shadedmathbox{E = mc^2} and \shadedmathbox[smugreen]{a^2 + b^2 = c^2}
    
    \vspace{0.5em}
    
    Display math: \shadedmathdisplay[scisyellow]{\int_0^\infty e^{-x^2} dx = \frac{\sqrt{\pi}}{2}}
\end{frame}

% ================================================================================
\section{Components Design}
% ================================================================================

\subsection{Block Environments}

\begin{frame}{Block Showcase}
    \begin{columns}[T]
        \begin{column}{0.48\textwidth}
            \begin{block}{Standard}
                This is a standard block with SMU styling.
            \end{block}
            
            \vspace{0.3em}
            
            \begin{alertblock}{Alert}
                Warning block for important information.
            \end{alertblock}
            
            \vspace{0.3em}
            
            \begin{examples}
                Example block for demonstrating concepts.
            \end{examples}
        \end{column}
        
        \begin{column}{0.48\textwidth}
            \begin{lightbox}
                Light box for highlighting supplementary information.
            \end{lightbox}
            
            \vspace{0.3em}
            
            \begin{tagblock}{Tag}
                Tagged block for categorized content.
            \end{tagblock}
            
            \vspace{0.3em}
            
            \begin{theorem}[Pythagorean Theorem]
                For a right triangle with sides $a$ and $b$, and hypotenuse $c$:
                \[ a^2 + b^2 = c^2 \]
            \end{theorem}
        \end{column}
    \end{columns}
\end{frame}

\subsection{Column Layouts}

\begin{frame}{Two Column Layout - Custom Width}
    The \smucode{\textbackslash begin\{twocolumns\}} command defaults to 50\% width for each column. You can use \smucode{\textbackslash leftcol[0.6]} and \smucode{\textbackslash rightcol[0.35]} to customize the column widths.
    
    \begin{twocolumns}
        \leftcol[0.6]{
            \begin{block}{Wider Left Column (60\%)}
                This column takes 60\% of the available width, leaving 35\% for the right column (plus some spacing).
                
                \vspace{0.5em}
                
                Useful for:
                \begin{itemize}
                    \item Main content vs. sidebar
                    \item Text vs. small figure
                    \item Detailed explanation vs. summary
                \end{itemize}
            \end{block}
        }
        \rightcol[0.35]{
            \begin{block}{Narrow Right (35\%)}
                Smaller column for supplementary information.
            \end{block}
        }
    \end{twocolumns}
\end{frame}

\begin{frame}{Three Column Layout}
    \begin{threecolumns}
        \leftcol[0.31]{
            \textbf{Step 1}
            
            Initialize
            
            \vspace{0.3em}
            
            \redbf{Input:} Data
        }
        \midcol[0.31]{
            \textbf{Step 2}
            
            Process
            
            \vspace{0.3em}
            
            \bluebf{Transform:} Apply
        }
        \rightcol[0.31]{
            \textbf{Step 3}
            
            Output
            
            \vspace{0.3em}
            
            \smucolor{smugreen}{\textbf{Result:} Done}
        }
    \end{threecolumns}
    
    \vspace{0.5em}
    
    \begin{lightbox}
        Three columns are perfect for showing processes or comparisons.
    \end{lightbox}
\end{frame}

% ================================================================================
\subsection{Tables and Lists}
% ================================================================================

\begin{frame}{Tables}
    \textbf{Example: Performance Comparison}
    
    \vspace{0.5em}
    
    \begin{table}
        \centering
        \begin{tabular}{lccc}
            \toprule
            \textbf{Method} & \textbf{Accuracy (\%)} & \textbf{Time (ms)} & \textbf{Memory (MB)} \\
            \midrule
            Method A & 95.3 & 12.5 & 128 \\
            Method B & 97.1 & 18.2 & 256 \\
            Method C & 96.8 & 15.7 & 192 \\
            \bluebf{Ours} & \bluebf{98.5} & \bluebf{14.3} & \bluebf{160} \\
            \bottomrule
        \end{tabular}
        \caption{Performance comparison of different methods}
    \end{table}
    
    \vspace{0.5em}
    
    \begin{lightbox}
        Use \texttt{booktabs} package for professional-looking tables.
    \end{lightbox}
\end{frame}

\begin{frame}{Mixed Lists in Columns}
    \begin{twocolumns}
        \leftcol{
            \textbf{Unordered Lists (Itemize):}
            \begin{itemize}
                \item First level item
                \item Another first level item
                \begin{itemize}
                    \item Second level nested item
                    \item Another second level item
                    \begin{itemize}
                        \item Third level deeply nested
                        \item More third level content
                    \end{itemize}
                    \item Back to second level
                \end{itemize}
                \item Back to first level
                \item Final first level item
            \end{itemize}
        }
        \rightcol{
            \textbf{Ordered Lists (Enumerate):}
            \begin{enumerate}
                \item First step in the process
                \item Second step with details
                \begin{enumerate}
                    \item Sub-step 2.1
                    \item Sub-step 2.2
                    \begin{enumerate}
                        \item Sub-sub-step 2.2.1
                        \item Sub-sub-step 2.2.2
                    \end{enumerate}
                \end{enumerate}
                \item Third step returning to main level
                \item Fourth and final step
            \end{enumerate}
        }
    \end{twocolumns}
\end{frame}

\begin{frame}{Description Lists}
    \textbf{Custom Description Lists:}
    
    \begin{description}
        \item[\bluebf{API}] Application Programming Interface for software communication
        \item[\bluebf{SDK}] Software Development Kit containing tools and libraries
        \item[\bluebf{IDE}] Integrated Development Environment for coding
        \item[\bluebf{CI/CD}] Continuous Integration and Continuous Deployment
    \end{description}
    
    \vspace{1em}
    
    \begin{lightbox}
        Description lists are perfect for glossaries or term definitions.
    \end{lightbox}
\end{frame}

% ================================================================================
\section{Code Rendering}
% ================================================================================

\subsection{Code Highlighting}

\begin{frame}[fragile]{Code Highlighting Features}
    The theme uses \textbf{minted} for professional code highlighting:
    
    \vspace{0.5em}
    
    \begin{twocolumns}
        \leftcol{
            \textbf{Features:}
            \begin{itemize}
                \item Line numbers
                \item Syntax highlighting
                \item 49 color themes
                \item 100+ languages
                \item Inline code support
            \end{itemize}
        }
        \rightcol{
            \textbf{Supported Languages:}
            \begin{itemize}
                \item Python, C, C++, Java
                \item JavaScript, TypeScript
                \item Go, Rust, Shell
                \item SQL, R, MATLAB
                \item And many more!
            \end{itemize}
        }
    \end{twocolumns}
    
    \vspace{0.5em}
    
    \begin{alertblock}{Important}
        Always use \texttt{[fragile]} option when frames contain code blocks!
    \end{alertblock}
\end{frame}

\begin{frame}[fragile]{Python: Fibonacci Sequence}
\begin{minted}{python}
def fibonacci(n):
    """Calculate the nth Fibonacci number using recursion"""
    if n <= 1:
        return n
    return fibonacci(n-1) + fibonacci(n-2)

def fibonacci_iterative(n):
    """Calculate Fibonacci iteratively (more efficient)"""
    a, b = 0, 1
    for _ in range(n):
        a, b = b, a + b
    return a

# Test the functions
for i in range(10):
    result = fibonacci_iterative(i)
    print("F(%d) = %d" % (i, result))
\end{minted}

Inline Python: \pyinline{lambda x: x**2}
\end{frame}

\begin{frame}[fragile]{C++: Template Function}
\begin{minted}{cpp}
#include <iostream>
#include <vector>
#include <algorithm>

template<typename T>
T findMax(const std::vector<T>& arr) {
    if (arr.empty()) {
        throw std::runtime_error("Empty vector");
    }
    return *std::max_element(arr.begin(), arr.end());
}

int main() {
    std::vector<int> numbers = {3, 7, 2, 9, 1, 5};
    std::cout << "Maximum: " << findMax(numbers) << std::endl;
    
    std::vector<double> values = {3.14, 2.71, 1.41};
    std::cout << "Maximum: " << findMax(values) << std::endl;
    
    return 0;
}
\end{minted}
\end{frame}

\begin{frame}[fragile]{SQL and Shell Scripts}
    \begin{columns}[T]
        \begin{column}{0.48\textwidth}
            \textbf{SQL Query}
\begin{minted}{sql}
SELECT 
    s.student_name,
    AVG(g.score) as avg_score,
    COUNT(*) as num_courses
FROM students s
JOIN grades g 
    ON s.id = g.student_id
WHERE g.score >= 80
GROUP BY s.id, s.student_name
HAVING COUNT(*) >= 3
ORDER BY avg_score DESC
LIMIT 10;
\end{minted}
        \end{column}
        
        \begin{column}{0.48\textwidth}
            \textbf{Bash Script}
\begin{minted}{bash}
#!/bin/bash

BACKUP_DIR="/backup"
DATE=$(date +%Y%m%d)
SOURCE="/data"

# Create backup
echo "Starting backup..."
tar -czf \
  "$BACKUP_DIR/bk_$DATE.tar.gz" \
  "$SOURCE"

if [ $? -eq 0 ]; then
    echo "Backup completed"
else
    echo "Backup failed!"
    exit 1
fi
\end{minted}
        \end{column}
    \end{columns}
\end{frame}

\begin{frame}[fragile]{Inline Code Examples}
    \textbf{Inline code for different languages:}
    
    \begin{itemize}
        \item Python: \pyinline{def hello(): return "world"}
        \item C: \cinline{int *ptr = malloc(sizeof(int));}
        \item Java: \javainline{List<String> names = new ArrayList<>();}
        \item JavaScript: \jsinline{const arr = [1, 2, 3].map(x => x * 2);}
    \end{itemize}
    
    \vspace{1em}
    
    \begin{block}{Generic Inline Code}
        You can also use \smucode{\textbackslash code[language]\{code\}} for any language:
        \begin{itemize}
            \item Generic: Use \code[C]{printf("Hello")} for output
            \item Commands: Run \smucode{git commit -m "message"}
            \item Paths: See file \smucode{/etc/config.conf}
        \end{itemize}
    \end{block}
\end{frame}

% ================================================================================
\subsection{Summary}
% ================================================================================

\begin{frame}{Theme Components Summary}
    \begin{twocolumns}
        \leftcol{
            \textbf{Text \& Colors:}
            \begin{itemize}
                \item SMU brand colors
                \item Custom color commands
                \item Text highlighting
                \item Bold color variants
                \item Math highlighting
            \end{itemize}
            
            \vspace{0.5em}
            
            \textbf{Blocks:}
            \begin{itemize}
                \item Standard blocks
                \item Alert blocks
                \item Example blocks
                \item Light boxes
                \item Tagged blocks
                \item Theorem boxes
            \end{itemize}
        }
        \rightcol{
            \textbf{Code:}
            \begin{itemize}
                \item Minted integration
                \item 49 color themes
                \item 100+ languages
                \item Line numbers
                \item Inline code
            \end{itemize}
            
            \vspace{0.5em}
            
            \textbf{Layout:}
            \begin{itemize}
                \item 2-column layouts
                \item 3-column layouts
                \item Custom widths
                \item Progress bar
                \item Section pages
            \end{itemize}
        }
    \end{twocolumns}
\end{frame}

\begin{frame}{Quick Reference}
    \begin{tagblock}{Getting Started}
        Load the theme: \smucode{\textbackslash usepackage\{smu\}}
    \end{tagblock}
    
    \vspace{0.3em}
    
    \begin{tagblock}{Code Highlighting}
        For code blocks: \smucode{\textbackslash begin\{frame\}[fragile]}
        
        Use environments: \smucode{pycode}, \smucode{cppcode}, \smucode{jscode}, etc.
    \end{tagblock}
    
    \vspace{0.3em}
    
    \begin{tagblock}{Columns}
        Two columns: \smucode{\textbackslash begin\{twocolumns\}...\textbackslash end\{twocolumns\}}
        
        Three columns: \smucode{\textbackslash begin\{threecolumns\}...\textbackslash end\{threecolumns\}}
    \end{tagblock}
    
    \vspace{0.3em}
    
    \begin{tagblock}{Documentation}
        See \texttt{README.md} for detailed documentation and examples.
    \end{tagblock}
\end{frame}

\begin{frame}{Customization}
    \textbf{The theme is highly customizable:}
    
    \vspace{0.5em}
    
    \begin{description}
        \item[\bluebf{Colors}] Edit \texttt{smucore/smucolors.def}
        \item[\bluebf{Fonts}] Edit \texttt{smucore/smufonts.def}
        \item[\bluebf{Code Style}] Edit \texttt{smucore/smucode.def}
        \item[\bluebf{Commands}] Edit \texttt{smucore/smucommands.def}
        \item[\bluebf{Blocks}] Edit \texttt{smucore/smublocks.def}
        \item[\bluebf{Templates}] Edit \texttt{smucore/smutemplates.def}
    \end{description}
    
    \vspace{1em}
    
    \begin{lightbox}
        All components are modular and can be customized independently!
    \end{lightbox}
\end{frame}

\begin{frame}{Resources}
    \begin{tagblock}{Documentation}
        \begin{itemize}
            \item \texttt{README.md} - Main documentation
            \item \texttt{main\_test.tex} - This demo file
            \item Source files in \texttt{smucore/} directory
        \end{itemize}
    \end{tagblock}
    
    \vspace{0.5em}
    
    \begin{tagblock}{External Resources}
        \begin{itemize}
            \item Minted: \url{https://ctan.org/pkg/minted}
            \item Beamer: \url{https://ctan.org/pkg/beamer}
            \item Pygments Styles: \url{https://pygments.org/styles/}
        \end{itemize}
    \end{tagblock}
\end{frame}

\begin{frame}{Reference Test}
\textbf{MtA-based Approaches (Multiplicative-to-Additive)}
    \begin{itemize}
        \item Paillier Encryption: Requires strong-RSA assumption.
        \\
        \cite{GGN16/gennaro2016threshold, Lin17/lindell2021fast, GG18/gennaro2018fast, CGG20+/canetti2020uc, XLX24/xie2024direct}
        \item CL Encryption: Based on Class Groups.
        \\
        \cite{CLL+19/castagnos2019two, CCL+20/castagnos2020bandwidth, YCX21/yuen2021compact, DMA+21/deng2021promise}
        \item Oblivious Transfer: No extra assumptions and computationally efficient.
        \\
        \cite{DKLs18/doerner2018secure, DKLs19/doerner2019threshold, DKLs23/doerner2024threshold}
        \item Joye-Libert: Based on revisited Joye-Libert encryption.
        \\
        \cite{XAL+23/xue2023efficient}
    \end{itemize}

\textbf{Threshold CL Encryption}
    \begin{itemize}
        \item Decryption key of CL shared among $n$ parties.
        \\
        \cite{BDO23/braun2023secure, WMY23/wong2023real, WMC24/wong2024secure, TX25/tang2025robust, JTX25/jiang2025three}
    \end{itemize}

\textbf{Others}
\begin{itemize}
    \item Pseudorandom Correlation Generator \cite{ANO+22/abram2022low}, ...
\end{itemize}
    
\end{frame}

\begin{frame}[allowframebreaks]{References}
    \linespread{0.8}\selectfont
    \setlength{\itemsep}{0pt}
    \setlength{\parskip}{0pt}
    \printbibliography[heading=none]
\end{frame}

% ================================================================================
% Thank You Page
% ================================================================================
\backmatter

\end{document}
